\subsection{Lower subsolutions}

The Bregman obstruction motivates returning to ordinary Euclidean proximal
subproblems while enforcing order explicitly.

\begin{definition}[Lower certificate]
A vector $\ell\ge0$ is a lower certificate for $F_\rho^+$ if
\begin{equation}
 \kappa_{\rho,i}(\ell)\ge0
 \qquad\text{for every }i\in\supp(\ell).
 \label{eq:lower-certificate}
\end{equation}
\end{definition}

\begin{lemma}[$M$-matrix lower comparison; \ProvedStatus]
\label{lem:lower-comparison}
If~\eqref{eq:lower-certificate} holds, then
\begin{equation}
  \ell\le x_\rho^\star
  \qquad\text{coordinatewise}.
  \label{eq:lower-order}
\end{equation}
\end{lemma}

\begin{proof}[Proof idea]
Suppose $(\ell-x_\rho^\star)_i$ attained a positive maximum after the usual
positive diagonal scaling.  At such a coordinate $i$ we have $\ell_i>0$,
so~\eqref{eq:lower-certificate} applies.  The nonpositive off-diagonal entries
of $Q$, together with the KKT inequalities~\eqref{eq:kkt}, force
$(Q(\ell-x_\rho^\star))_i\le0$, contradicting positive definiteness and the
$M$-matrix maximum principle.  Equivalently, this is the standard comparison
argument for obstacle problems with a nonsingular $M$-matrix.
\end{proof}

\subsection{A local Euclidean proximal subproblem}

Given a lower certificate $\ell$, define
\begin{equation}
 z^\star
 :=\arg\min_{z\ge0}
 \left\{F_\rho(z)
 +\frac{\kappa_{\mathrm A}}{2}\norm{z-\ell}_2^2\right\}.
 \label{eq:safe-euclidean-prox}
\end{equation}

\begin{theorem}[Order-safe Euclidean center; \ProvedStatus]
\label{thm:safe-euclidean-center}
If $\ell$ satisfies~\eqref{eq:lower-certificate}, then
\begin{equation}
  \ell\le z^\star\le x_\rho^\star
  \label{eq:prox-sandwich}
\end{equation}
and
\begin{equation}
  p(z^\star)+\rho\vol(\supp z^\star)\le1.
  \label{eq:euclidean-prox-budget}
\end{equation}
\end{theorem}

\begin{proof}
At any coordinate where $z^\star$ would cross below $\ell$, the optimality
system of~\eqref{eq:safe-euclidean-prox}, the lower-certificate inequality,
and the $M$-matrix maximum principle give a contradiction.  Comparing the
same system with the KKT system of $x_\rho^\star$ gives the upper inequality in
~\eqref{eq:prox-sandwich}.

For the volume bound, write the proximal KKT equation as
\[
 (Q+\kappa_{\mathrm A} I)z^\star
 =b+\kappa_{\mathrm A}\ell-\alpha\rho\,w\odot u,
\]
where $u_i=1$ on $\supp(z^\star)$ and $u_i\in[0,1]$ elsewhere.  Multiplying by
$w^T$ gives
\[
 (\alpha+\kappa_{\mathrm A})p(z^\star)
 +\alpha\rho\sum_i d_i u_i
 =\alpha+\kappa_{\mathrm A} p(\ell).
\]
Since $p(z^\star)\ge p(\ell)$ by~\eqref{eq:prox-sandwich},
\[
 \alpha\left(p(z^\star)+\rho\sum_i d_i u_i\right)\le\alpha.
\]
This implies~\eqref{eq:euclidean-prox-budget}.
\end{proof}

Ordinary Euclidean proximal subproblems therefore already have the desired
$1/\rho$ locality whenever their centers remain lower certificates.  The
problem is acceleration: removing extrapolation gives the usual proximal
point contraction $1-\Theta(\alpha)$ and total work
$\softO(1/(\alpha\rho))$.

\subsection{Complete conditional variants}

Two conditional theorems survive all later counterexamples.

\begin{theorem}[Bounded-core warmup; \ConditionalStatus]
\label{thm:bounded-core}
Let $\bar\rho=(1+\sqrt\alpha)\rho$.  Suppose a known region $C$ contains
$\supp(x_{\bar\rho}^\star)$ and
\begin{equation}
  \vol(C)\le B/\rho.
  \label{eq:bounded-core}
\end{equation}
If every accelerated warmup subproblem is restricted to $C$, then a
fixed-face Euclidean accelerated method reaches~\eqref{eq:warm-gap} in
\begin{equation}
 \Work_{\mathrm{warm}}
 =\softO\!\left(\frac{B}{\rho\sqrt\alpha}\right).
 \label{eq:bounded-core-work}
\end{equation}
After a lower retraction, the full hybrid method has
\[
 \Work_{\mathrm{hybrid}}
 =\softO\!\left(\frac{B+1}{\rho\sqrt\alpha}\right).
\]
\end{theorem}

\begin{theorem}[Controlled face phases; \ConditionalStatus]
\label{thm:controlled-faces}
Suppose a certified active-set method reaches~\eqref{eq:warm-gap} after at
most $K$ face enlargements, every face is contained in
$\supp(x_{\bar\rho}^\star)$, and momentum is restarted at each enlargement.
Then
\begin{equation}
 \Work_{\mathrm{hybrid}}
 =\softO\!\left(\frac{K+1}{\rho\sqrt\alpha}\right).
 \label{eq:face-phase-work}
\end{equation}
\end{theorem}

Both statements are fixed-face acceleration plus the support
budget~\eqref{eq:rppr-support-budget} and the tail
bound~\eqref{eq:tail-target}.  They are consistent with expanding-subspace
acceleration \citep{martinezrubio2023accelerated}.  Neither theorem bounds
$B$ or $K$ uniformly over graphs.
